% ============================================================================
%  ANEXO 1.A — PLAN INSTITUCIONAL DE CERTIFICACIÓN ISO/IEC 27001:2022
%  audIT Soluciones Tecnológicas SpA - Licitación TFEP-01/2026
% ============================================================================

\section*{Anexo 1.A: Plan institucional de certificación ISO/IEC 27001:2022}\label{sec:1-anexo-iso}
\addcontentsline{toc}{section}{Anexo 1.A: Plan institucional de certificación ISO/IEC 27001:2022}

Conforme a la exigencia de acreditación de Requisito Habilitante del \art{34.1}{22} de las Bases Administrativas TFEP-01/2026, audIT formaliza el siguiente plan institucional vinculante.

\subsection*{1. Individualización del oferente y representación legal}
El plan identifica al proponente, a su representante y al organismo a cargo del proceso de certificación.
\begin{itemize}
  \item \textbf{Razón Social:} audIT Soluciones Tecnológicas SpA
  \item \textbf{RUT:} 76.924.310-0
  \item \textbf{Representante Legal:} Alejandro Hermosilla Díaz
  \item \textbf{Cédula de Identidad:} 14.892.341-8
  \item \textbf{Organismo Certificador:} Bureau Veritas Certification S.A., organismo a cargo del proceso institucional de certificación.
  \item \textbf{Norma Internacional:} ISO/IEC 27001:2022 (\textit{Information Security, Cybersecurity and Privacy Protection - Information Security Management Systems}).
  \item \textbf{Expediente de Auditoría Externa:} N.° BV-EXP-2026-CL-8921.
\end{itemize}

\subsection*{2. Alcance del Sistema de Gestión de Seguridad de la Información (SGSI)}
El Sistema de Gestión de Seguridad de la Información de audIT cubre:
\begin{quote}
«El diseño, desarrollo, pruebas, integración, aseguramiento de calidad, implantación y operación de plataformas de software de misión crítica, soluciones de telemetría e Internet de las Cosas (IoT) de borde, y administración de arquitecturas cloud para transporte terrestre, distribución y logística corporativa, conforme a la Declaración de Aplicabilidad (SoA) versión 3.0».
\end{quote}

\subsection*{3. Estado documental del proceso y fundamentación del Plan (\art{34.1}{22})}
audIT acoge formalmente su acreditación a la disposición expresa del \art{34.1}{22} de las Bases Administrativas, que faculta la presentación de un plan institucional de certificación con hitos verificables dentro de los primeros doce meses del Contrato. A la fecha de presentación de esta propuesta técnica, la compañía acredita el siguiente estado documental y procedimental ante Bureau Veritas Certification S.A. (Expediente de Auditoría N.° BV-EXP-2026-CL-8921):
\begin{itemize}
  \item \textbf{Auditoría de Fase 1 (Revisión Documental y Diseño del SGSI):} Concluida y aprobada conforme en noviembre de 2025.
  \item \textbf{Auditoría de Fase 2:} En curso ante Bureau Veritas Certification S.A. La evaluación, el cierre de observaciones y la recomendación del equipo auditor forman parte del expediente; no se declara la certificación como emitida.
  \item \textbf{Estado del Expediente y Compromiso de Emisión:} El certificado oficial no se encuentra emitido a la fecha de la propuesta. audIT compromete su entrega en el \textbf{Mes 1 de Contrato}, como obligación de su plan institucional y no como una fecha impuesta por el artículo 34.1. El expediente BV-EXP-2026-CL-8921 y la declaración del representante se individualizan en el Sobre N.º 1.
\end{itemize}

\subsection*{4. Cronograma vinculante de despliegue, protocolización y vigilancia anual}
Para dar cabal cumplimiento a lo dispuesto en el \art{34.1}{22} de las Bases Administrativas (\textit{«plan de certificación con hitos verificables dentro de los primeros doce meses del Contrato»}), en la \tabref{1-iso27001-hitos} se formaliza el cronograma de cumplimiento de los cuatro hitos comprometidos.

\begin{tabla}[ancho=completo]{Cronograma vinculante de despliegue, protocolización y vigilancia anual ISO/IEC 27001:2022}{1-iso27001-hitos}{C{15mm} L{32mm} X L{35mm} X}
\textbf{Hito} & \textbf{Plazo Máximo} & \textbf{Actividad y Entregable Verificable} & \textbf{Responsable Corporativo} & \textbf{Criterio de Verificación / Evidencia} \\ \midrule
\textbf{H1} & Mes 1 de Contrato & \textbf{Protocolización y Entrega de Certificado Oficial:} Recepción del certificado emitido por Bureau Veritas, con alcance, vigencia e identificador verificable, entregando copia autorizada a Transportes Curimón S.A. & Representante Legal / Oficial de Seguridad (CISO) & Certificado formal oficial con identificador verificable del organismo certificador. \\
\textbf{H2} & Mes 4 de Contrato & \textbf{Auditoría Interna de Extensión al Proyecto Curimón:} Evaluación formal de los 93 controles ISO 27001 aplicados a los componentes de nube, borde y datos de la solución Curimón. & Líder de Calidad y Seguridad de audIT SpA & Informe pericial de auditoría interna de proyecto con matriz de riesgos mitigada. \\
\textbf{H3} & Mes 8 de Contrato & \textbf{Revisión por la Dirección y Actualización de SoA:} Evaluación de métricas operacionales de seguridad, análisis de vulnerabilidades OWASP ASVS y actualización de Declaración de Aplicabilidad. & Comité de Seguridad de la Información (CSIC) & Acta formal firmada por Gerencia General y CISO, con matriz de tratamiento de riesgos actualizada. \\
\textbf{H4} & Mes 11 de Contrato & \textbf{Auditoría Externa Anual de Mantenimiento:} Ejecución de la primera auditoría periódica de vigilancia por parte de Bureau Veritas Certification S.A. para certificar la mejora continua. & Bureau Veritas Certification S.A. / CISO & Informe de Auditoría Externa de Vigilancia Anual favorable y sin observaciones bloqueantes. \\
\end{tabla}

Como se desprende de la \tabref{1-iso27001-hitos}, la emisión se compromete en M1 y los controles de extensión y vigilancia se realizan en M4, M8 y M11. Los identificadores H1 a H4 de esta tabla son hitos del plan institucional ISO 27001; no sustituyen los hitos de aceptación o pago del contrato.

\subsection*{5. Declaración formal y firma del representante legal}
En mi calidad de Representante Legal de audIT, declaro bajo fe de juramento que la información expuesta precedentemente es fidedigna y que la empresa asume la obligación de dar cumplimiento a los hitos del presente Plan dentro de los plazos señalados, facultando a Transportes Curimón S.A. para auditar su avance en cualquier momento durante la vigencia del Contrato.

Santiago de Chile, 03 de octubre de 2026.

\vspace{0.8cm}
\begin{center}
  \includegraphics[height=1.5cm]{portadas/activos/media-firma.png}\\[4pt]
  \textbf{Alejandro Hermosilla Díaz}\\
  Representante Legal y Director Ejecutivo\\
  audIT Soluciones Tecnológicas SpA\\
  RUT: 14.892.341-8
\end{center}
